\section{5b. Bases de Datos Documentales}

\begin{frame}{5.1 Concepto: El documento autodescriptivo}
    A diferencia de SQL, aquí no guardamos ``filas'', guardamos \textbf{Documentos} (normalmente JSON o BSON).
    \begin{itemize}
        \item \textbf{Autodescriptivo:} Cada documento lleva sus propias ``columnas'' (llaves) dentro.
        \item \textbf{Jerárquico:} Puedes meter listas y otros objetos dentro de un documento.
        \item \textbf{SGBD Referente:} \textbf{MongoDB}.
    \end{itemize}
    \begin{exampleblock}{La analogía del Archivador}
        SQL es una hoja de Excel rígida. NoSQL Documental es una carpeta con folios: cada folio puede tener una estructura distinta, pero todos están guardados en el mismo cajón.
    \end{exampleblock}
\end{frame}

\begin{frame}{5.2 Esquema Flexible (Schema-less)}
    ¿Qué significa realmente que no hay esquema?
    \begin{itemize}
        \item \textbf{Polimorfismo:} En la colección \texttt{alumnos}, un alumno puede tener \texttt{telefono} y otro puede tener un array \texttt{telefonos: [``casa'', ``movil'']}. A la base de datos le da igual.
        \item \textbf{Evolución sin dolor:} Si tu App necesita un campo nuevo, lo añades en el código y listo. No hay que hacer un \texttt{ALTER TABLE} que bloquee la base de datos durante horas.
    \end{itemize}
    \begin{alertblock}{La responsabilidad es del programador}
        Como la base de datos no valida el esquema, es tu código el que debe asegurarse de que los datos tienen sentido.
    \end{alertblock}
\end{frame}

\begin{frame}{5.3 Modelado: ¿Embeber o Referenciar?}
    Esta es la pregunta del millón en NoSQL Documental:
    \begin{itemize}
        \item \textbf{Embeber (Embedding):} Metes los datos hijos dentro del padre.
        \begin{itemize}
            \item \textit{Cuándo:} Relaciones 1:1 o 1:pocos (ej: direcciones de un usuario).
            \item \textit{Ventaja:} Lees todo de un golpe (Sin JOINs).
        \end{itemize}
        \item \textbf{Referenciar (Linking):} Guardas el ID del hijo.
        \begin{itemize}
            \item \textit{Cuándo:} Relaciones 1:Muchos que crecen sin parar (ej: comentarios de un post viral).
            \item \textit{Ventaja:} Evitas que el documento supere el límite de tamaño (16MB en Mongo).
        \end{itemize}
    \end{itemize}
\end{frame}

\begin{frame}{5.4 Indexación y Potencia de Consulta}
    Aunque sea NoSQL, puedes buscar de forma muy avanzada:
    \begin{itemize}
        \item \textbf{Índices secundarios:} Puedes crear un índice sobre el email, la edad o incluso sobre un campo que está dentro de una lista.
        \item \textbf{Agregaciones:} Puedes filtrar, agrupar y calcular medias como en SQL, pero usando un ``pipeline'' (tubería) de etapas.
    \end{itemize}
    \begin{exampleblock}{Consultas Ad-hoc}
        A diferencia de Clave-Valor, aquí puedes decir: ``Dame todos los usuarios de Madrid que tengan más de 20 años y les guste el Rock''.
    \end{exampleblock}
\end{frame}

\begin{frame}[fragile]{5.5 Ejemplo Práctico (MongoDB Shell)}
    \begin{lstlisting}[language=SQL]
// Insertar un profesor con sus modulos embebidos
db.profesores.insertOne({
  nombre: "Vicente",
  especialidad: "Informatica",
  modulos: [
    { nombre: "Bases de Datos", horas: 6 },
    { nombre: "Sostenibilidad", horas: 3 }
  ],
  activo: true
})

// Buscar profesores que den mas de 5 horas de Bases de Datos
db.profesores.find({
  "modulos.nombre": "Bases de Datos",
  "modulos.horas": { $gt: 5 }
})
    \end{lstlisting}
\end{frame}
